% Part: sets-functions-relations
% Chapter: arithmetization
% Section: cauchy

\documentclass[../../../include/open-logic-section]{subfiles}

\begin{document}

\olfileid{sfr}{arith}{cauchy}
\olsection{Bijlage: reële getallen opbouwen met Cauchyrijen}
In \olref[cuts]{sec} bouwden we de reële getallen op met
Dedekindsneden. Hier leggen we een andere manier uit. We beginnen bij
de definitie van Cauchy van wat we nu een Cauchyrij noemen. Andere
schrijvers uit de negentiende eeuw gebruikten die definitie om de reële
getallen te \emph{bouwen}. Dat waren vooral Weierstrass, Heine,
M\'{e}ray en Cantor. (Een goed overzicht van deze geschiedenis staat
bij \citeauthor{OConnorRobertson:RN}
\citeyear{OConnorRobertson:RN}.)
Maar eerst gaan we terug naar Simon Stevin (1548--1620). Stevin zag dat
we elk reëel getal kunnen weergeven met de cijfers van zijn
decimaalontwikkeling. Zelfs een irrationaal getal als $\sqrt{2}$ heeft
zo'n rij cijfers. Die begint zo:
\[
	1.41421356237\ldots
\] 
Zo'n rij cijfers past heel makkelijk in de verzamelingenleer. Zie hem
als een functie $d \colon \Nat \to \Nat$. Dan is $d(n)$ het cijfer op
de $n$-de plaats achter de komma waarnaar we kijken. We moeten nog wel
iets toevoegen voor het stuk vóór de komma, hier alleen $1$. Ook
hebben we een equivalentierelatie nodig: een regel die ervoor zorgt dat
verschillende schrijfwijzen voor hetzelfde getal bij elkaar horen. Zo
moet bijvoorbeeld $0.999\ldots = 1$ gelden. Op deze manier kunnen we de
reële getallen in de verzamelingenleer volledig nauwkeurig opbouwen,
op Stevins manier.
Stevin is hier niet ons hoofdonderwerp. (Meer over Stevin staat in
\citealt{KatzKatz2012}.) Maar zijn idee brengt ons dicht bij een andere
mogelijkheid. In plaats van meteen de cijfers van $\sqrt{2}$ te
gebruiken, nemen we een \emph{rij}: een lijst in een vaste volgorde van
steeds nauwkeurigere rationale benaderingen van $\sqrt{2}$. We schrijven
telkens een groter stuk uit:
\[
1, 1.4, 1.414, 1.4142, 1.41421,\ldots
\]
Het idee van ``steeds betere benaderingen'' geeft een paar nieuwe
stappen. Ze lijken op de stappen waarmee we $\Rat$ uit $\Int$ en
$\Int$ uit $\Nat$ bouwden. Dit zijn de belangrijkste:
\begin{enumerate}
	\item Elk reëel getal kun je schrijven als een decimaalontwikkeling.
	Die kan een oneindige rij cijfers hebben.
	\item Dezelfde informatie kun je vastleggen met een rij rationale
	getallen.
	\item Zo'n rij rationale getallen kun je zien als een functie van
	$\Nat$ naar $\Rat$. Neem gewoon $f(n)$ als het $n$-de rationale getal
	in de rij.
\end{enumerate}
Maar niet \emph{elke} functie van $\Nat$ naar $\Rat$ geeft een reëel
getal. Neem bijvoorbeeld deze functie:
\[
	f(n) =\begin{cases}
			1 & \text{als }n\text{ oneven is}\\
			0 &\text{als }n\text{ even is}
		\end{cases}
\]
De rij $0,1,0,1,0,1,0,\ldots$ komt niet steeds dichter bij één reëel
getal. We willen alleen rijen gebruiken die dat wel doen. Daarom kijken
we alleen naar rijen met een \emph{limiet}: een getal waar de termen
uiteindelijk zo dicht bij komen als je maar wilt.
In \olref[his][set][limits]{sec} kwamen we een limiet al tegen. Toch
kunnen we hier niet \emph{helemaal} dezelfde definitie gebruiken. In
``$(\forall \epsilon>0)$'' liep de variabele stilzwijgend over de reële
getallen. Ook keken we daar naar functies op reële getallen. Als we die
definitie nu gebruiken, nemen we de reële getallen dus al aan. Maar die
proberen we juist te \emph{bouwen}. Gelukkig kunnen we een bijna gelijk
idee gebruiken.
\begin{defn}\ollabel{def:CauchySequence}
  Een functie $f: \Nat \to \Rat$ heet een \emph{Cauchyrij} precies als
  voor elk positief getal $\epsilon \in \Rat$ geldt: $(\exists \ell
  \in \Nat)(\forall m, n > \ell)|f(m) - f(n)| < \epsilon$.
\end{defn}
Het idee blijft hetzelfde. Kies hoe nauwkeurig je wilt zijn; dat meten
we met~$\epsilon$. Daarna is er een ``gebied'' waarin je moet kijken:
alle invoeren groter dan~$\ell$. Onze rij $1$, $1.4$, $1.414$,
$1.4142$, $1.41421$\ldots heeft duidelijk een limiet. Wil je
$\sqrt{2}$ benaderen met een fout van minder dan
$\nicefrac{1}{10^{n}}$, dan is elke term na de $n$-de nauwkeurig genoeg.
Je zou nu kunnen zeggen dat elk reëel getal gewoon een Cauchyrij
\emph{is}. Maar net als bij onze opbouw van $\Int$ en $\Rat$ is dat te
na\"ief. Verschillende Cauchyrijen kunnen hetzelfde reële getal
aangeven. Dat zie je zo. Neem een Cauchyrij~$f$. Maak $g$ precies gelijk
aan~$f$, behalve dat $g(0)\neq f(0)$. Na het eerste getal zijn de twee
rijen overal gelijk. Uiteindelijk willen we dus zeggen dat ze dezelfde
limiet hebben, in de betekenis van \olref{def:CauchySequence}. Ze
``definiëren'' dan hetzelfde reële getal. Daarom moeten we die twee
Cauchyrijen als hetzelfde reële getal zien.
Daarom moeten we weer een equivalentierelatie maken voor de
Cauchyrijen: een regel die zegt welke rijen bij elkaar horen. Daarna
stellen we reële getallen gelijk aan klassen van Cauchyrijen die volgens
die regel bij elkaar horen. Eerst moeten
we zeggen wanneer een functie uiteindelijk naar $0$ gaat. Neem een
functie $h : \Nat \to \Rat$. We zeggen dat \emph{$h$ naar $0$ gaat}
precies als voor elk positief getal $\epsilon \in \Rat$ geldt:
$(\exists \ell \in \Nat)(\forall n > \ell)|h(n)| <
\epsilon$.\footnote{Vergelijk dit met de definitie van $\lim_{x
\mathord{\rightarrow}\infty}f(x) = 0$ in
\olref[his][set][limits]{sec}.} Neem verder functies $f$ en $g$ van
$\Nat \to \Rat$ en stel $(f-g)(n) = f(n) - g(n)$. Nu leggen we vast:
\[
	f \Realequiv g \text{ precies als $(f-g)$ naar $0$ gaat}.
\]
We moeten nog nagaan dat $\Realequiv$ echt een equivalentierelatie is.
Dat klopt. Daarna kunnen we de reële getallen zien als de klassen van
Cauchyrijen die volgens $\Realequiv$ bij elkaar horen, voor alle
Cauchyrijen van $\Nat \to \Rat$.
\begin{prob}
Neem $f(n) = 0$ voor elke $n$. Neem $g(n) = \frac{1}{(n+1)^2}$. Laat
zien dat beide functies Cauchyrijen zijn en allebei limiet $0$ hebben.
Daarom geldt ook $f \sim_\Real g$.
\end{prob}
Zoals altijd schrijven we $\equivrep{f}{\Realequiv}$ voor de klasse
waarin $f$ !!a{element} is. Om de tekst leesbaar te houden, laten we de
index hierna weg. We schrijven dus alleen $\equivrep{f}{}$. Ook spreken
we af dat voor elke $q \in \Rat$ geldt: $q_{\Real} =
\equivrep{c_{q}}{}$. Hier is $c_{q}$ de constante functie waarvoor
$c_q(n) = q$ voor elke $n \in \Nat$. Daarna definiëren we de gewone
relaties en rekenbewerkingen voor reële getallen, bijvoorbeeld:
\begin{align*}
	\equivrep{f}{} + \equivrep{g}{} &= 	\equivrep{(f + g)}{} \\
	\equivrep{f}{} \times 	\equivrep{g}{} &= \equivrep{(f \times g)}{} 
\end{align*}
Hierbij is $(f + g)(n) = f(n) + g(n)$ en $(f \times g)(n) = f(n)
\times g(n)$. We moeten ook nagaan dat $(f + g)$, $(f-g)$ en $(f\times
g)$ weer Cauchyrijen zijn als $f$ en $g$ dat zijn. Dat klopt; het
bewijs laten we als oefening.
Tot slot leggen we vast wanneer het ene getal kleiner is dan het andere. We noemen $\equivrep{f}{}$
\emph{positief} precies als $\equivrep{f}{}\neq 0_\Real$ én $(\exists
\ell \in \Nat)(\forall n > \ell)0 < f(n)$. Daarna zeggen we dat
$\equivrep{f}{} < \equivrep{g}{}$ precies als $\equivrep{(g - f)}{}$
positief is. We moeten nagaan dat deze regel altijd dezelfde uitkomst
geeft. Hij mag dus niet afhangen van welke functie uit de klasse we kiezen.
%\begin{align*}
%	\equivrep{f}{} \leq \equivrep{g}{} \text{ iff }&\equivrep{f}{} = \equivrep{g}{} \text{ or }\\
%	&(\exists \ell \in \Nat)(\forall n \in \Nat)(\ell < n \lif f(n) \leq g(n))
%\end{align*}
Als dat klaar is, kunnen we vrij makkelijk laten zien dat de gewone
rekenregels gelden:
\begin{thm}\ollabel{thm:cauchyorderedfield}
De Cauchyrijen vormen een geordend lichaam.
\end{thm}
\begin{proof}
Oefening.
\end{proof}
\begin{prob}
Bewijs dat de Cauchyrijen een geordend lichaam vormen.
\end{prob}
Het is lastiger om te bewijzen dat de reële getallen die we zo hebben
gebouwd de volledigheidseigenschap hebben. Daarom geven we dat bewijs.
\begin{thm}
Elke niet-lege verzameling Cauchyrijen die een bovengrens heeft, heeft
ook een kleinste bovengrens.
\end{thm}
\begin{proof}[Schets van het bewijs] Neem een niet-lege verzameling
Cauchyrijen $S$ met een bovengrens. Er is dus een $p \in \Rat$ waarvoor
$p_{\Real}$ een bovengrens van $S$ is. Neem $r \in S$; dan is er een
$q \in \Rat$ waarvoor $q_{\Real} < r$. Als $S$ een kleinste bovengrens
heeft, ligt die dus tussen $q_\Real$ en $p_\Real$. De twee grenspunten
tellen mee.
We komen steeds dichter bij de kleinste bovengrens door haar tegelijk
van onderen en van boven te benaderen. We maken daarvoor twee functies,
$f, g \colon \Nat \to \Rat$. De functie $f$ moet de kleinste
bovengrens\ van bovenaf benaderen en $g$ van onderaf. We beginnen met:
\begin{align*}
	f(0) &= p \\
	g(0) &= q
\end{align*}
Neem daarna $a_n = \frac{f(n) + g(n)}{2}$ en leg het volgende
vast:\footnote{Dit is een recursieve definitie: elke volgende stap
gebruikt de vorige. Maar we hebben \emph{nog} niet laten zien waarom
zo'n definitie mag.}
\begin{align*}
	f(n+1) &=
	\begin{cases}
		a_n &\text{als }(\forall h \in S)\equivrep{h}{} \leq (a_n)_\Real\\
		f(n)&\text{anders}
	\end{cases}\\
	g(n+1) &=
	\begin{cases}
		a_n &\text{als }(\exists h \in S)\equivrep{h}{} \geq (a_n)_\Real\\
	 	g(n) &\text{anders}
	\end{cases}
\end{align*}
Zowel $f$ als $g$ is een Cauchyrij. (Dat kun je vrij makkelijk
nagaan; we laten het als oefening.) De functie $(f-g)$ gaat naar $0$,
want het verschil tussen $f$ en $g$ is na elke stap hoogstens half zo groot.
Daarom is $\equivrep{f}{} = \equivrep{g}{}$.
Eerst laten we zien dat $\equivrep{f}{}$ een bovengrens van $S$ is.
Dat betekent\ dat $(\forall h \in S)\equivrep{h}{} \leq
\equivrep{f}{}$. (Onderweg gebruiken we
\olref{thm:cauchyorderedfield}.) Neem $h \in S$ en stel, om een
tegenspraak te krijgen, dat $\equivrep{f}{} < \equivrep{h}{}$. Dan is
$0_\Real < \equivrep{(h-f)}{}$. De Cauchyrij $f$ blijft gelijk of
daalt. Daarom is er een $n \in \Nat$ waarvoor
$\equivrep{(c_{f(n)} - f)}{} < \equivrep{(h-f)}{}$. Dus:
\[
	(f(n))_\Real = \equivrep{c_{f(n)}}{} < \equivrep{f}{} + \equivrep{(h-f)}{} = \equivrep{h}{},
\]
maar volgens onze bouw geldt juist $\equivrep{h}{} \leq
(f(n))_\Real$. Dat is een tegenspraak.
Nu laten we zien dat $\equivrep{f}{} = \equivrep{g}{}$ de
\emph{kleinste} bovengrens van $S$ is. Neem een willekeurige Cauchyrij
$j$ en stel dat $\equivrep{j}{} < \equivrep{g}{}$. Net als hierboven
volgt, omdat $g$ \emph{gelijk blijft of stijgt}, dat er een $n \in \Nat$ is
waarvoor $\equivrep{j}{} < (g(n))_\Real$. Maar door onze bouw is er een
$h \in S$ waarvoor $(g(n))_\Real \leq \equivrep{h}{}$. Dus
$\equivrep{j}{} < \equivrep{h}{}$. Daarom is $\equivrep{j}{}$ geen
bovengrens van $S$.
\end{proof}
\end{document}

